% Pista geo-26j-q4b · Geometria
% Procedència: PAU Catalunya 2026 · Convocatòria de juny · Sèrie 1 · Exercici 4 — Opció B
\documentclass[11pt,a4paper]{article}
\usepackage{pista}
\pistainfo{Catalunya 2026}{Convocatòria de juny}{Sèrie 1}

\begin{document}

\section*{\textcolor{blauPista}{Exercici 4 --- Opció B}}

\noindent Considereu els punts de l'espai $P = (1, 0, -1)$, $Q = (3, -2, 0)$ i $R = (1, 1, 1)$.

\subsection*{4.1. \normalfont Calculeu l'equació del pla que conté els punts $P$, $Q$ i $R$. \hspace*{\fill}\textnormal{[0,75~punts]}}

\pista{Els vectors $\vec{PQ}$ i $\vec{PR}$ són vectors directors del pla. En primer lloc, calcula'ls. A continuació, escriu l'equació del pla que passa pel punt $P$ i que té aquests dos vectors directors. Per obtenir l'equació, iguala a zero el determinant:
\[
\begin{vmatrix} x - 1 & y - 0 & z + 1 \\ \ldots & \ldots & \ldots \\ \ldots & \ldots & \ldots \end{vmatrix} = 0
\]
Finalment, desenvolupa'l i escriu el pla en la forma \mbox{$Ax + By + Cz + D = 0$}.}

\subsection*{4.2. \normalfont Comproveu que l'àrea del triangle $\Delta PQR$ és $\dfrac{3\sqrt{5}}{2}\,u^{2}$. \hspace*{\fill}\textnormal{[0,75~punts]}}

\pista{Fes servir aquesta fórmula, amb els vectors $\vec{PQ}$ i $\vec{PR}$ de l'apartat 4.1:
\[
\text{Àrea} = \frac{1}{2}\left|\vec{PQ} \times \vec{PR}\right|
\]
En primer lloc, calcula el producte vectorial. A continuació, calcula'n el mòdul. Finalment multiplica'l per $\frac{1}{2}$.}

\subsection*{4.3. \normalfont Determineu les condicions que han de complir les coordenades d'un quart punt $S = (x, y, z)$ per tal que $P$, $Q$, $R$ i $S$ formin un tetraedre de volum 1. (El volum del tetraedre format pels punts $P$, $Q$, $R$ i $S$ és: $\text{volum}(PQRS) = \dfrac{\text{àrea}(\Delta PQR) \cdot \text{altura}}{3}$.) \hspace*{\fill}\textnormal{[1~punt]}}

\pista{A la fórmula que et donen, l'altura és la distància del punt $S$ al pla que conté $P$, $Q$ i $R$ (el de l'apartat 4.1).
\begin{itemize}
\item Iguala el volum a 1.
\item Substitueix-hi l'àrea de l'apartat 4.2.
\item Per a l'altura, escriu la distància punt-pla en funció de les coordenades $x$, $y$, $z$.
\end{itemize}
T'apareixerà una equació amb un valor absolut. Recorda que una equació com $|X| = 5$ vol dir que $X = 5$ o bé $X = -5$. Per tant, hauràs de resoldre dues equacions, i obtindràs dos plans.}

\end{document}
